% =============================================================================
%  Supplementary Section S8: Robustness and audit of the numbers.
% =============================================================================
\section{Robustness and audit of the numbers}
\label{si:s8}

Every computed number in the main text and in this Supplementary Information is printed by a
script or follows by stated arithmetic from printed numbers; literature values carry their
citations. This section describes how the numbers were produced (Section~\ref{si:s8-how}),
where the Fisher information was checked against simulated fits and why simulation was used
at 5\% of normal and for one-sided statements (Section~\ref{si:s8-fim}), which corrections
made during the audit changed a reported value (Section~\ref{si:s8-corr}), and which script
produces each headline number of the main text (Section~\ref{si:s8-table},
Supplementary Table~\ref{stab:s8-headline}). Section~\ref{si:s8-limits} lists what the audit
did not check.

% -----------------------------------------------------------------------------
\subsection{How the numbers were produced}
\label{si:s8-how}

\paragraph{Simulated records.} No patient data enter the analysis. A record is one
realisation of the recorded load of a single simulated patient, with a 10\% day-to-day
variation about its mean, sampled on a stated schedule over two years. Twenty-four such
intake records (random seeds 0--23) are used, all generated with the same constants
($k_e=5$\,d$^{-1}$, $K_m=100$\uM), so they describe one patient's possible diets, not 24
patients. Unless stated otherwise, a precision is the median over the 24 records, with its
range. For values from the Fisher information each record is scored at the true parameters;
nothing is fitted.
% src: audit/brief_2026-10-01/REPORT.md:139 (24 draws of one simulated patient's intake, all k_e = 5; no patient data), :143 (each seed scored by its own FIM at the truth; nothing fitted), :1367 (no patient data); paper/sections/02_model.tex:371-373 (median over 24 intake records with range)

\paragraph{Two layers of scripts.} In the original analysis, the model and design code
(\nolinkurl{src/common}, \nolinkurl{src/model_A}, \nolinkurl{src/model_C}, \nolinkurl{src/design})
and the analysis scripts in \nolinkurl{audit/} write the numbers its text quotes into
\nolinkurl{results/numbers/*.tex}, each file naming the script that wrote it (one,
\nolinkurl{drift.tex}, records its generator as unregistered). The constants of the
linear-clearance and carnitine models live in one registry (\nolinkurl{src/common/params.py}),
and tests assert that its values are the ones the simulation code runs on; the two constants
of the saturable-clearance model, $V_e$ and $K_e$, are defaults in
\nolinkurl{src/common/models.py} with no row in that registry (Supplementary
Table~\ref{stab:params}). The audit (\nolinkurl{audit/brief_2026-10-01}) re-derives and
extends that analysis in 60 scripts (\nolinkurl{t0_footing.py} to
\nolinkurl{t56_contrast_mc.py}), with 13 further recheck scripts; each prints its results to
\nolinkurl{out/<script>.txt}, and every number in the audit report appears in one of those
files. Where the audit and the original analysis disagree, the audit's value is used. The
main text and this Supplementary Information transcribe their numbers from these files;
Supplementary Table~\ref{stab:s8-headline} gives the file and line behind each headline
number.
% src: results/numbers/*.tex headers ('generated by ... -- do not edit'; drift.tex: 'an unregistered generator'); src/common/params.py:59-61 (tests), :66 (imports ParamsA, ParamsC only); REPORT.md:774 (Model-B parameter set unregistered); audit/brief_2026-10-01/README.md:5,98; REPORT.md:3; ls audit/brief_2026-10-01/t*.py: 60 scripts plus 13 *_recheck*

\paragraph{The Fisher-information footing.} Precision is computed from the Fisher information
of the steady-state balance in the parameters $(\ln\theta,\ln K_m,b,\ln k_e)$, with $b$ the
offset that carries the unrecorded load. The observation noise combines the 5\% assay
coefficient of variation with the 5\% error in the recorded load; $b$ has a normal prior with
mean zero and standard deviation 5\% of the load, $\ln k_e$ a normal prior centred on its
true value with standard deviation 0.8, and $\theta$ and $K_m$ no prior. The audit's implementation reproduces the original analysis to $10^{-13}$, and a
finite-difference version written on $\ln m$ agrees to $3.9\times10^{-8}$. Replacing the
single pooled noise level by a noise level for each draw gives 3.045 instead of 3.067 at 5\%
of normal, which is the difference between the 3.05 and 3.07 quoted in the main text. Every
record carries the known fortnightly dose train ($D=6000$\um, half-life 3.5\,d); without it
the value at 5\% of normal would be 8.59, and with a constant load as well the information
matrix is singular.
% src: REPORT.md:3 (core.fim reproduces n4.sd_ln_theta to 1e-13), :140-141 (flux-space FIM over (ln th, ln Km, b, ln ke); b prior 5% L; K_m no prior; 3.9e-8; 3.045 vs 3.067; without dose train 8.59; singular); out/t1_identifiability.txt:18-19 (8.594; per-draw 3.0454); src/common/params.py:110-111 (KE_PRIOR_LOG 0.8, LOAD_FRAC 0.05)

\paragraph{Numerical methods.} Trajectories are integrated with adaptive stiff solvers
(\texttt{LSODA} or \texttt{Radau}) at a relative tolerance of $10^{-8}$ or tighter. Two
regression tests assert that halving the integration step moves an answer by less than
$10^{-6}$ relative and that an integrated steady state equals the closed-form steady state;
these check the arithmetic and say nothing about the model's validity. Random draws use \nolinkurl{numpy.random.default_rng}
with an explicit seed at each call site. Critical capacities are found by root-finding on the
reduced carnitine clearance (Supplementary Section~\ref{si:s1-carn}), with $R_0$ re-derived
at every setting of the shape constants; at $\theta=\thci$ the peak of the clearance equals
the load to ten digits.
% src: paper/sections/02_model.tex:348-374; paper/sections/03_dynamic.tex:286-292 (peak of C at theta_crit_int equals the load to ten digits); out/t53_crossing_shape.txt:8 (R0 re-derived each time)

\paragraph{Independent re-derivation.} The audit's follow-up scripts on estimators, the
cohort and the Fisher information (\nolinkurl{t6c}, \nolinkurl{t6d}, \nolinkurl{t8}--\nolinkurl{t11})
and its third-round scripts (\nolinkurl{t13}--\nolinkurl{t19}) were re-derived by independent
verifiers along routes sharing no code with them (the \nolinkurl{*_recheck} scripts), as were
the scripts on the in vitro to in vivo mapping, the breath test and the saturable-clearance
model (\nolinkurl{t29}, \nolinkurl{t31}, \nolinkurl{t33}, \nolinkurl{t35}, \nolinkurl{t36}) and on the
carnitine model's two-observable information (\nolinkurl{t41}). The comparison with the
published model structures (\nolinkurl{audit/k2_structures.py} and
\nolinkurl{audit/k5_identifiability.py}) was rechecked by \nolinkurl{audit/k3_recheck.py} and
\nolinkurl{audit/k6_recheck.py}. The Monte Carlo crossing for a direct flux reading
(\nolinkurl{t12}) has no independent recheck; no main-text number depends on it.
% src: REPORT.md:3 (follow-up and round-3 scripts re-derived; t12 not), :1404 (t12 no recheck), :1410 (round-3 *_recheck3), :1418 (t29, t31 verified), :1419 (t33/t35/t36 verified), :1423 (t41 verified); audit/PHASE4_PROPOSED_TEXT.md:318-319 and audit/PHASE5_ROBUSTNESS.md:151-152 (k2, k5 rechecked by k3, k6)

Supplementary Table~\ref{stab:s8-methods} summarises which method produced each kind of number.

\begin{table}[tbp]
\centering
\footnotesize
\caption{\label{stab:s8-methods}How each kind of number was computed. FIM: Fisher information
matrix; CRB: Cram\'er--Rao bound, the inverse of the FIM; NLS: nonlinear least squares; MAP:
maximum a posteriori fit.}
\begin{tabular}{@{}>{\raggedright\arraybackslash}p{0.21\linewidth}
>{\raggedright\arraybackslash}p{0.27\linewidth}
>{\raggedright\arraybackslash}p{0.23\linewidth}
>{\raggedright\arraybackslash}p{0.21\linewidth}@{}}
\toprule
kind of number & method & main-text examples & check \\
\midrule
log-scale precision, $\sdl{\hat\theta}$, $\sdl{\hat\lambda}$ & CRB on the log scale, median
over 24 records & sampling rules; 82--627\% band; 3.07 at 5\%; measurement routes; $\lambda$
from the record & per-draw noise model; NLS spot-check; line geometry \\
\addlinespace
one-sided bound on $\theta$ & profile likelihood, 24 records (22\%); the balance, the
non-negativity of $k_e$ and $K_m$ from the record (30.7\%) & capacity below 22\% of normal
(30\% without the $k_e$ prior) & noiseless expected profile, also 22.4\% \\
\addlinespace
within-patient change & FIM, then 400 simulated record pairs per scenario fitted by
constrained MAP & 20\um; 56\um\ at 80\% power; false-alarm rates & difference against log
ratio \\
\addlinespace
one patient's 21-day decision & NLS on 200 simulated records per point & estimating $\theta$
does not rescue 21 days & delta method and log-normal propagation compared \\
\addlinespace
crossing half-lives, critical capacities & closed form for the trough; root-finding on the
reduced clearance; $R_0$ re-derived at each corner & Table~1 of the main text & peak of
clearance equals load at $\thci$ \\
\addlinespace
late-infusion peaks & integration of the two-state carnitine model; delay to doubling by
bisection to 0.025\,d & how late an infusion can be & $k=1$ and $k=10$; validity limit \\
\addlinespace
separation of model structures & records simulated from the two published structures and
fitted with our model; one post-infusion draw by integration & $\hat\theta$ 1.3--2.0 times
too high at our dose, 4--20 at the calibrated dose; 16.3 standard deviations at 3 minutes &
independent recheck scripts \\
\addlinespace
probability of clearing a crossing & log-normal distribution for the hepatic half-life about
4.3\,d, standard deviation of the log 0.35--0.53 & 64--71\% & a calculation on a literature
band; no data enter it \\
\bottomrule
\end{tabular}
\end{table}
% src: REPORT.md:7 (how to read the numbers), :972-998 (R-M); out/t1_identifiability.txt:54-57, :68 (expected-profile upper end 672, 22.4%); out/t18_set_id.txt:9-12; out/t56_contrast_mc.txt:4-13; out/t46_timing.txt:24-31; out/t53_crossing_shape.txt:5-6,80; out/t55_infusion_tolerance.txt:4-8; results/numbers/structure.tex:9; paper/sections/A1_supplementary.tex:73-75

% -----------------------------------------------------------------------------
\subsection{Fisher information against simulated fits}
\label{si:s8-fim}

\paragraph{The rule.} The Cram\'er--Rao bound (CRB), the inverse of the Fisher information,
describes a symmetric, locally Gaussian uncertainty. For a standard deviation on the log
scale it is reliable in this problem, except where the estimate can reach zero (next
paragraph), and shrinks roughly as $1/\sqrt n$ with the number of draws $n$. It is not reliable for a quantity on the natural scale, or for a one-sided
statement built on the same posterior, once the log-scale standard deviation is large, as at
5\% of normal, where $\sdl{\hat\theta}\approx3$. The posterior is then close to log-normal,
with a long right tail on the natural scale and a left tail that reaches zero, so a lower
bound is set by the tails and does not tighten as the log-scale width shrinks. Such
quantities were read from simulated fits (Monte Carlo). Where an estimate can sit on the
boundary $\theta=0$, the usual asymptotic theory of likelihood intervals also fails, which is
a second reason to simulate~\cite{boundary_lrt,boundary_bootstrap}.
% src: REPORT.md:7 (log-scale sd reliable except where theta_hat can reach 0, F6), :976 (the lesson, stated generally); manuscript/main.tex:696-701

\paragraph{Where the CRB is optimistic.} Given $K_m$ and $k_e$, $\theta$ enters the
steady-state balance linearly, so the Fisher information describes a Gaussian likelihood in
$\theta$ itself, and simulated fits attain it there (the ratio of root-mean-square error to
the bound was 0.99--1.06 across the scenarios checked). At 5\% of normal that Gaussian is
wide compared with $\theta$, so estimates at or below zero are common, and any summary on the
log scale depends on how the constraint $\theta\ge0$ is handled. Even with $k_e$ known
exactly or to 5\%, the profile likelihood reaches $\theta$ near zero in 28--38\% of records,
so the symmetric log-scale interval of the CRB is optimistic on the low side.
% src: REPORT.md:405-407 (theta enters linearly; FIM describes a Gaussian in linear theta), :423 (linear RMSE/FIM 0.99-1.06; profile reaches theta near 0 in 28-38% of records; FIM interval optimistic on the low side)

At 5\% of normal on the monthly record, the CRB of 3.067 would imply a symmetric 95\%
interval whose upper end is 166,704 times its lower end (a factor of about 408 on each side
of the estimate). The profile likelihood over the 24
records instead has its 95\% upper end at a median of 672\um\ (22.4\% of normal; range
451--821\um) and reaches the lowest value of the grid ($\theta=0.05$\um) in all 24 records.
The maximum-likelihood estimate has a median of 3.8\um, and 54\% of records put it below a
tenth of the truth. The main text therefore reports only a one-sided bound.
Without the prior on $k_e$, the balance, the non-negativity of $k_e$ and the value of $K_m$
recovered from the record bound $\theta$ at 30.7\% of normal.
% src: out/t1_identifiability.txt:52-57 (CRB 3.067, factor 166,704 = exp(2 x 1.96 x 3.067), upper end over lower end, t1_identifiability.py:174; exp(1.96 x 3.067) = 408; upper end median 672 = 22.4%, range 451-821; grid floor 24/24; MLE median 3.8, 54% below 15); out/t18_set_id.txt:5-12 (30.7% non-negativity; 22.4% with k_e prior)

Nonlinear least-squares fits to 300 simulated two-year monthly records at 5\% of normal,
with the half-life also estimated, gave a realised $\sdl{\hat\theta}$ of 8.26 against the
Fisher value of about 3.05. The median of $\hat\theta/\theta$ was 1.08, its 5--95\% range
0.00--2.97, and only 8.7\% of fits fell within $\pm20\%$ of the truth. Simulation therefore
supports the negative conclusion of the main text. Freeing the half-life inflates the 8.26,
so the 8.26 does not revise the 3.05, which is computed with the half-life known. The
pairwise intersections of the line geometry (Supplementary Section~\ref{si:s1-pencil}) give
the same picture without any noise model.
% src: out/t47_fim_audit.txt:16-22 (300 records, 24 months monthly, median 1.08, 5-95% [0.00, 2.97], 8.7% within +/-20%, sd 8.26); REPORT.md:992, :1429 (8.3 is lambda-free and overstates the lambda-fixed 3.05; not a revision)

\paragraph{One patient's decision on the 21-day interval.} Whether one patient's estimate
supports a 21-day interval turns on the lower bound of the trough margin
$\theta+\dtr(21)-\thci$, a one-sided quantity on the natural scale. From the same
information, a delta-method lower bound was about $-440$\um\ and a log-normal propagation
gave about $+19$\um, a disagreement of about 450\um. Nonlinear least-squares fits to 200
simulated records per point settled it in favour of the more pessimistic reading: the standard
deviation of the estimated margin stays between 109 and 162\um\ from six months to five years
of weekly draws, and its 5th percentile stays below zero at every capacity (conditions and
results in Supplementary Section~\ref{si:s6:theta}). This is the basis of the main-text
statement that estimating $\theta$ does not rescue the 21-day interval.
% src: REPORT.md:976 (delta-method -440 vs log-normal MC +19, ~450 uM/d), :1430 (t46: NLS, 200 records per point), :939-970 (R-L: lambda estimated from the same record); out/t46_timing.txt:12-22 (theta_base 150/65/20, weekly, 6-60 months; printed sd(mu_hat) 109-162; p5 < 0; the reading at :24-25 rounds this to ~105-165)

\paragraph{Within-patient change.} Simulated record pairs (Supplementary
Section~\ref{si:s5:pairs}) confirm the Fisher value for the difference
$\hat\theta_2-\hat\theta_1$ (a standard deviation of 19.4--21.3\um\ against
$0.133\times150=20$\um) but not for the log ratio, whose false-alarm rate is 34\% rather
than 5\% (Supplementary Section~\ref{si:s5:ratio}). The change is therefore reported in \um,
and the main text's drift and variability figures are the Fisher values read as multiples of
150\um.
% src: out/t56_contrast_mc.txt:4-13 (0.133, 20.0 uM/d; table: floor 20-28%, mean diff +1.0/+31.2/-30.3/+66.9/-46.0 against 0/+30/-30/+67.5/-45, sd 19.4-21.3, reject ln 0.34 diff 0.05 at no change; tails to about +-10, line 8); REPORT.md:1194-1206 (R-U)

\paragraph{Hidden correlations.} The posterior of the five-parameter on-therapy problem at 5\%
of normal (24 months of weekly draws, half-life estimated) has two strong correlations
involving $\theta$, $-0.99$ between $\ln\theta$ and $\ln k_e$ and $+0.87$ between $\ln\theta$
and $\ln\lambda$; the others that are strong ($-0.85$ between $\ln k_e$ and $\ln\lambda$,
$+0.88$ between $\ln K_m$ and $b$) do not feed a statement about $\theta$. The first feeds the
margin $\theta-\theta_{\mathrm{crit}}$, which the original analysis simulated, and the second
feeds the 21-day margin above, so no strong correlation feeds a natural-scale statement
propagated only through the Fisher information. In the linear-clearance model, a classification
of patients between 7-, 10- and 14-day intervals (critical capacity 430\um, capacities spread
log-uniformly over 5--150\um, 24 months of monthly draws) built on log-normal draws with the
Fisher width gave 61\% correct and none unambiguous; redone with nonlinear least-squares
estimates on simulated records it gave 54\% and none, so that pessimistic result survives.
% src: out/t47_fim_audit.txt:28-46 (correlation matrix at theta = 150, 24 mo weekly: -0.99, +0.87, k_e-lambda -0.85, Km-b +0.88; theta-k_e feeds the paper's simulated mu = theta - theta_crit, theta-lambda feeds mu_long), :5-13 (spot-check A: grid {7,10,14}, theta_crit 430, theta_base log-uniform 5-150, 24 mo monthly; 54.3% vs 61%, 0%); REPORT.md:990-996

\paragraph{What remains Fisher-based.} Two one-sided figures have not been simulated: the
chance that one patient's record places the half-life on one side of a crossing
(Supplementary Section~\ref{si:s6:record}), and the number of pooled patients needed to do so
(23--155 in the main text, depending on cadence, record length and power, and only if all
share one half-life). They are indicative. The probabilities of clearing a crossing
(64--71\% and the variants in the main text) are calculations on an assumed log-normal band
for the hepatic half-life; no data enter them.
% src: REPORT.md:1435 (t52 record-resolution numbers are FIM sds, not simulated; side-of-crossing chances one-sided, not simulated); out/t52_knife_resolution.txt:66-68 (77/155 monthly 12 mo, 35/71 weekly 12 mo, 23/46 weekly 24 mo); out/t53_crossing_shape.txt:80-81,89-91

% -----------------------------------------------------------------------------
\subsection{Audit corrections that changed reported values}
\label{si:s8-corr}

The analysis was audited in successive rounds, each re-deriving or extending earlier results.
Several reported values changed, in both directions: some earlier statements were too
favourable and others too pessimistic. Supplementary Table~\ref{stab:s8-corrections} lists the
corrections that bear on numbers in the main text or in this Supplementary Information, with
the cause of each and the script that now produces the value. The main text and this
Supplementary Information use the corrected values throughout.
% src: REPORT.md:96-135 (what changed across the rounds), :133 (reversals in both directions); REPORT.md:1167-1206 (R-T, R-U)

\begingroup
\footnotesize
\setlength{\LTcapwidth}{\linewidth}
\renewcommand{\arraystretch}{1.1}
\begin{longtable}{@{}>{\raggedright\arraybackslash}p{0.23\linewidth}
>{\raggedright\arraybackslash}p{0.30\linewidth}
>{\raggedright\arraybackslash}p{0.24\linewidth}
>{\raggedright\arraybackslash}p{0.12\linewidth}@{}}
\caption{\label{stab:s8-corrections}Audit corrections that changed a reported value. Scripts
are those of the audit unless a path is given.}\\
\toprule
earlier statement & corrected value & cause & script \\
\midrule
\endfirsthead
\multicolumn{4}{@{}l}{\emph{Supplementary Table~\thetable, continued}}\\
\toprule
earlier statement & corrected value & cause & script \\
\midrule
\endhead
\midrule
\multicolumn{4}{r@{}}{\emph{continued on the next page}}\\
\endfoot
\bottomrule
\endlastfoot
within-patient change detectable at 80\% power: a rise of 45\% or a fall of 30\% &
$\pm56$\um\ (37\% of a patient at 5\% of normal); a 45\% rise is detected 91\% and a 30\% fall
64\% of the time & a Fisher log-ratio value read as a ratio; simulated, the change is measured
as a difference & \nolinkurl{t56} \\
an unrecorded antibiotic course mimics a 3- to 20-fold change in capacity & an apparent rise
of one to three times the patient's whole capacity (a linear-response estimate, for a 10--41\%
fall in the load) & the same misreading, a linearised Fisher shift of 1.1 to 2.9 log units
read as a ratio & \nolinkurl{t24}, \nolinkurl{t56} \\
within-patient precision worse under saturable clearance (0.133 to 0.167) & better, 0.088 &
the load offset was not shared between the two windows (a coding error) & \nolinkurl{t33},
\nolinkurl{t37} \\
21-day interval deliverable to a growing subgroup as the record lengthens (51--73\%) &
not supportable from the passive record at any length & a one-sided margin propagated through
the Fisher information; NLS simulation & \nolinkurl{t45}, \nolinkurl{t46} \\
PCC half-life assumed from rodent pharmacokinetics, scaled allometrically & measured in human
fibroblasts, 3.6\,d~\cite{chandler1985}; protein degradation rates transfer without
allometric scaling & literature search & \nolinkurl{t48} \\
hepatic PCC likely turns over faster than fibroblast PCC & equal or slower, central estimate
about 4.3\,d (3--6\,d) & mitochondrial-matrix comparator, 4.26\,d~\cite{kim2012}, in place of
the whole-liver median & \nolinkurl{t50} \\
the on-therapy record settles which side of the 21-day crossing a patient is on & per
patient, a true half-life of 3.6\,d is placed 0.16--0.58 standard deviations from the
crossing over the records examined (0.58 after five years of weekly draws); pooled, 23--77
patients at 50\% power or 46--155 at 80\%, if all share one half-life & the cited 7--18
standard deviations separated 3.6 from 7\,d, a gap 43 times larger in log distance &
\nolinkurl{t52} \\
3.6\,d clears the 21-day crossing by 1.5\% & only at the nominal constants; over the assumed
ranges of the carnitine constants the crossing spans 3.29--3.95\,d & the carnitine shape
constants were not propagated & \nolinkurl{t53} \\
$\thci$ varies within a factor of 1.4 over the shape constants (12 random draws) &
109.9--232.1\um, a factor of 2.1, at the corners of the assumed ranges & too few draws &
\nolinkurl{t53} \\
capped dose 2672\um\ scaled to 15 and 21 days at the same rate & ceiling applied at each
interval: crossings 3.30, 3.54 and 4.95\,d at the nominal constants & the scaled doses broke
the ceiling & \nolinkurl{t53} \\
an on-time infusion absorbs an illness (peak $1.0\times$ usual) & peak 2.4--4.1 times usual at
the assumed dose and nominal constants (0--5\% of normal, half-life 3.5 or 4.3\,d, $k=1$) &
the peak window closed 3\,d after the dose and missed the rise at the end of the illness; a
doubling bisection at 0.2\,d tolerance printed four delays 0.1\,d low & \nolinkurl{t55} \\
an ex vivo PCC assay at 20\% repeatability reaches the target & only if the scale from cell
activity to whole-body capacity is known; otherwise 3.05; scaled to healthy controls, 0.65 &
the 20\% is repeatability only; the per-patient in vitro to in vivo scatter does not cancel &
\nolinkurl{t26}, \nolinkurl{t29} \\
breath-test repeatability 12.8\% & $14.02\pm12.86\%$ (mean and SD of per-subject
coefficients of variation, between visits, methylmalonic acidaemia)~\cite{breath_test_mma} &
12.8\% was the SD, not the mean & \nolinkurl{audit/t2_tier1_margin.py} \\
interval comparisons at a fixed jump per infusion & same delivered amount per day ($D\propto\tau$),
which reproduces the 87\% trough cut at 28 days (88\% from the closed form) & convention of
the original analysis & \nolinkurl{t44} \\
upper bound on capacity 24\% of normal (708\um, median of 6 records) & 22.4\% (672\um,
median of 24 records) & more records & \nolinkurl{t1} \\
best precision 0.123 at 193\% of normal & 0.123 at 212\% & a 22-point grid replaced by a bounded
minimisation; the minimum is flat & \nolinkurl{t1} \\
the linear-clearance model is the most favourable case for $\theta$ & true for the joint
conditioning; the marginal precision on $\theta$ improves under saturable clearance (0.29--0.92
against 3.05), because the enzyme's share of the load rises (0.16 to 0.37) & two different
objects & \nolinkurl{t33}, \nolinkurl{t36} \\
\end{longtable}
\endgroup
% src: REPORT.md:100-131 (changed-claims table rows: ex vivo :112, ratio route :114, Model-A favourable :117, Model-B contrast :118, fixed amplitude :123, 21-day subgroup :126, lambda :127, hepatic direction :129, a fortiori :130, 1.5% :131), :133 (reversals in both directions); REPORT.md:1131-1147 (R-R; z 0.16-0.58 at 3.6 d in out/t52_knife_resolution.txt:13-30), :1149-1165 (R-S incl. ceiling addendum), :1167-1192 (R-T correction: 2.4-4.1x at the booked dose and nominal constants), :1194-1206 (R-U), :144 (n6 708 = 24% on 6 records), :188 (minimum 0.1226 at 212%); out/t24_contrast_stress.txt:42-45 (antibiotic depth 10-41%, +1.14 to +2.94); out/t55_infusion_tolerance.txt:64-71 (ill on time 2.4-4.1x); out/t1_identifiability.txt:17,55; results/numbers/scope.tex:9-10 (band_best 0.123 at 193%); audit/n4_theta_scope.py:348 (22-point scan); out/t53_crossing_shape.txt:54,76-77,89-91; out/t56_contrast_mc.txt:8-12,23-24; out/t26_multiobservable.txt:5,8; out/t29_mapping_scatter.txt:18; out/t44_interval_decision.txt:6; out/t33_model_b.txt:41,58; out/t35_managed_load.txt:48; paper/sections/08_specification.tex:130-134 (14.02 +/- 12.86; 12.8 is the SD); results/numbers/tier1.tex:1 (generated by audit/t2_tier1_margin.py)

% -----------------------------------------------------------------------------
\subsection{Main-text numbers and the scripts that produce them}
\label{si:s8-table}

Supplementary Table~\ref{stab:s8-headline} lists the headline numbers of the main text with the
conditions under which each holds, the method and the file that prints it, with line
numbers. \nolinkurl{out/tN} stands for the printed output of audit script \nolinkurl{tN}, in
\nolinkurl{audit/brief_2026-10-01/out}; \nolinkurl{design.tex}, \nolinkurl{scope.tex},
\nolinkurl{tier1.tex} and \nolinkurl{structure.tex} are in \nolinkurl{results/numbers},
each written by the generator it names. Unless stated, the
patient is at the stated fraction of normal capacity in the linear-clearance model, on the
assumed fortnightly dose train with a 3.5-day half-life, with the measurement model and priors
of Supplementary Table~\ref{stab:params}.

\begingroup
\footnotesize
\setlength{\LTcapwidth}{\linewidth}
\renewcommand{\arraystretch}{1.1}
\begin{longtable}{@{}>{\raggedright\arraybackslash}p{0.26\linewidth}
>{\raggedright\arraybackslash}p{0.13\linewidth}
>{\raggedright\arraybackslash}p{0.21\linewidth}
>{\raggedright\arraybackslash}p{0.07\linewidth}
>{\raggedright\arraybackslash}p{0.21\linewidth}@{}}
\caption{\label{stab:s8-headline}Headline numbers of the main text and their sources. Method:
FIM, Fisher information (Cram\'er--Rao bound); PL, profile likelihood; MC, simulated fits;
ODE, integration of the model; CF, closed form; RF, root-finding.}\\
\toprule
statement & value & conditions & method & printed in \\
\midrule
\endfirsthead
\multicolumn{5}{@{}l}{\emph{Supplementary Table~\thetable, continued}}\\
\toprule
statement & value & conditions & method & printed in \\
\midrule
\endhead
\midrule
\multicolumn{5}{r@{}}{\emph{continued on the next page}}\\
\endfoot
\bottomrule
\endlastfoot
\multicolumn{5}{@{}l}{\emph{The infusions are the identification experiment}}\\
infusions against the marker's own movement & 42 times & normal capacity, monthly & FIM &
\nolinkurl{results/numbers/design.tex}:19 \\
draws at the dose interval against monthly & 27 times less precise (21--33) & normal capacity
& FIM & \texttt{design.tex}:5 \\
bookings every 14.5\,d, 51 visits & 0.16 (0.161--0.162) & normal capacity, two years & FIM &
\texttt{design.tex}:52, 59 \\
monthly draws, 25 visits & 0.17 (0.168--0.171) & normal capacity & FIM & \texttt{design.tex}:88 \\
widest $k_e$ prior that still meets the target & factor 2.67 & normal capacity, monthly & FIM &
\texttt{design.tex}:42 \\
visits to the target with the increment known & 11 & normal capacity, monthly & FIM &
\nolinkurl{results/numbers/tier1.tex}:12 \\
accuracy needed on the delivered increment & 12\% (11.4--11.9) & normal capacity, monthly &
FIM & \texttt{design.tex}:35--37 \\
monthly draws needed at a 14\% increment error & 118 (9.8 years) & normal capacity, monthly &
FIM & \texttt{tier1.tex}:11, 13 \\
intake accuracy needed & about 10\% & normal capacity, monthly & FIM & \texttt{design.tex}:49 \\
$K_m$ recovered by the record & 6--7\% & every capacity & FIM & \nolinkurl{out/t18_set_id.txt}:23; \nolinkurl{paper/sections/00_abstract.tex}:26 \\
28-day interval at the same delivered rate & precision $+6.5\%$, trough $-87\%$ & normal capacity
& FIM, CF & \nolinkurl{paper/sections/07_design.tex}:324--330; \nolinkurl{out/t44}:6 \\
\addlinespace
\multicolumn{5}{@{}l}{\emph{Published structures}}\\
ranking of schedules & unchanged & both published structures, 12 visits & MC &
\nolinkurl{audit/PHASE5_ROBUSTNESS.md}:592--593 (from \nolinkurl{audit/k2_structures.py}) \\
monthly and phase-spread designs & within about 3\% of the linear-clearance model &
algebraic structure, 60\% calibration, 12 visits; not the turnover structure & MC &
\nolinkurl{audit/PHASE5_ROBUSTNESS.md}:581--583, 593--594 \\
$\hat\theta$ under the algebraic structure & 1.3--2.0 times too high at our dose; 4--20 at
the calibrated dose & our priors on $k_e$ and the load & MC & \nolinkurl{audit/k2_structures.py},
\nolinkurl{k5_identifiability.py} \\
one draw after an infusion separates the families & 16.3 SD (9.7 at 15\,min, 0.68 at 1\,h) &
draw at 3\,min, 5\% assay noise & ODE & \nolinkurl{results/numbers/structure.tex}:5--9 \\
\addlinespace
\multicolumn{5}{@{}l}{\emph{Residual capacity in severe PA}}\\
band in which the target is met & 82--627\% of normal & monthly & FIM & \nolinkurl{out/t1}:17;
\texttt{scope.tex}:14, 17 \\
best precision & 0.123 at 212\% & monthly & FIM & \nolinkurl{out/t1}:17 \\
precision at 5\% of normal & 3.07 & monthly & FIM & \nolinkurl{out/t1}:25; \texttt{scope.tex}:23 \\
fifteenfold denser sampling & lower edge 82 to 75\% & every 2\,d, 365 draws & FIM &
\nolinkurl{out/t3_variance.txt}:34 \\
correlation of $\ln\hat\theta$ with $\ln\hat k_e$ & $-0.99$ & 5\% & FIM & \nolinkurl{out/t1}:25 \\
one-sided bound & 22\% (30\% without the $k_e$ prior) & 5\%, 95\% & PL; CF &
\nolinkurl{out/t1}:55; \nolinkurl{out/t18}:9--12 \\
enzyme's share of the load & 16\% & 5\%, untreated & CF & \nolinkurl{out/t1}:25;
\texttt{scope.tex}:25 \\
\addlinespace
\multicolumn{5}{@{}l}{\emph{Measurement routes, 5\% of normal, monthly}}\\
$k_e$ tracer at 5.4\%; load tracer at 10\%; both; $k_e$ and unrecorded load exact & 0.45;
3.04; 0.44; 0.29 & & FIM &
\nolinkurl{out/t26}:10, 12, 16, 17 \\
ex vivo assay at 20\% with its scale known & 0.20 & & FIM & \nolinkurl{out/t26}:8 \\
assay scaled to healthy controls & 0.65 (0.50--0.80) & mapping $R^2=0.327$; between-patient
sd 0.6--1.0; 30 controls & FIM & \nolinkurl{out/t29}:18 \\
saturable-clearance model & 0.29 (constructed point); 0.92 (nominal) & $K_e$ fixed & FIM &
\nolinkurl{out/t35}:48; \nolinkurl{out/t33}:41 \\
\addlinespace
\multicolumn{5}{@{}l}{\emph{Within-patient change, two windows of 104 weekly draws, 5\%}}\\
standard deviation of the difference & 20\um\ (0.133) & & FIM, MC & \nolinkurl{out/t56}:4, 8 \\
change detected at 80\% power & 56\um & two-sided 5\% & MC & \nolinkurl{out/t56}:23--24 \\
false alarms with no change: difference; log ratio & 5\%; 34\% & 400 pairs & MC &
\nolinkurl{out/t56}:8 \\
at a 26.4\% per-draw marker CV & 71\um\ (0.475) & CV of serum propionylcarnitine within a
day in healthy adults & FIM & \nolinkurl{out/t42}:2, 10 \\
$k_e$ drift 0.06; calibration shift 1.5\% & 42\um; 30\um & & FIM & \nolinkurl{out/t30}:13, 23 \\
unrecorded antibiotic course & apparent rise of one to three times capacity & load 10--41\%
lower in the second window; linear response & FIM & \nolinkurl{out/t24}:42--45;
\nolinkurl{REPORT.md}:1206 \\
saturable-clearance model & 0.088 & nominal constants & FIM & \nolinkurl{out/t33}:58 \\
\addlinespace
\multicolumn{5}{@{}l}{\emph{Dosing interval, carnitine model}}\\
intrinsic critical capacity & 150.8\um & baseline load 600\um, nominal constants & RF &
\nolinkurl{out/t53}:5 \\
crossing half-lives at 14, 15, 21\,d & 2.62, 2.75, 3.55\,d & assumed dose, same delivered rate
& CF & \nolinkurl{out/t53}:6 \\
21-day crossing; $\thci$ over the assumed ranges & 3.29--3.95\,d; 110--232\um & eight corners &
CF, RF & \nolinkurl{out/t53}:76--77 \\
crossings capped at the ceiling & 3.30, 3.54, 4.95\,d (corners 2.98--3.87, 3.19--4.15,
4.47--5.80) & ceiling at each interval & CF & \nolinkurl{out/t53}:89--91 \\
peak at the assumed dose & 2.18 times normal & 5\% of normal, 14\,d, half-life 3.5\,d & CF &
\nolinkurl{out/t44}:5, with $(150+6400)/3000$ \\
own capacity needed for 21 days at the capped dose & about 1.8\% of normal & half-life 4.3\,d,
nominal constants & CF & \nolinkurl{pi_note/pi_note.tex}:395--396 \\
chance the hepatic half-life clears the 21-day crossing & 64--71\% (56--60\% worst corner;
34--39\% capped) & log-normal about 4.3\,d, sd of log 0.35--0.53 & CF &
\nolinkurl{out/t53}:80--81, 91 \\
illness threshold & 994\um & linear-clearance model, $f=1.5$, $m_{\mathrm{tol}}=40$\uM & CF &
\nolinkurl{out/t44}:40 \\
crossings against 994\um & about 5.0, 5.2, 6.3\,d & assumed dose & CF &
\nolinkurl{pi_note/pi_note.tex}:430--433 \\
half-life from the record & sd of log 0.055 & 5\% of normal, two years, monthly, dated
events (jitter, missed doses, a supply gap) & FIM & \nolinkurl{out/t45}:47 \\
pooled patients to place a 3.6-day half-life & 23--155 & 50--80\% power at $2\sigma$; 12--24
months of monthly or weekly draws; one half-life shared by all & FIM &
\nolinkurl{out/t52}:66--68 \\
\addlinespace
\multicolumn{5}{@{}l}{\emph{Late infusions, two-state carnitine model, $k=1$, nominal constants}}\\
rise in the peak per day of delay & about 10\% & no own enzyme, half-life 4.3\,d, assumed
dose, no illness & ODE & \nolinkurl{out/t55}:179 \\
delay before the peak doubles & 6.5--8.8\,d & 0--5\% of normal, half-life 4.3\,d, assumed dose &
ODE & \nolinkurl{out/t55}:17--20 \\
skipped infusion & 5.1 to 2.7 times the usual peak & 0 to 5\% of normal, half-life 4.3\,d,
assumed dose & ODE & \nolinkurl{out/t55}:65--71 \\
illness: on time; 3\,d late; 7\,d late & 2.6, 4.7, 21.0 times the usual peak with no own
enzyme; 2.5, 4.1, 14.7 at 2.5\%; 2.4, 3.7, 10.3 at 5\% & half-life 4.3\,d, assumed dose,
illness 1.5 times the load for 14\,d & ODE & \nolinkurl{out/t55}:65--71 \\
least favourable setting & doubling after 2.4\,d (4.3\,d at a 4.3-day half-life) & $D=2672$\um,
worst corner, 3.5\,d; no own enzyme & ODE & \nolinkurl{out/t55}:46, 142, 153 \\
trough test fails before the peak doubles & in 6 of 8 settings (all but the assumed dose at
4.3\,d) & no own enzyme; assumed or capped dose, nominal or least favourable constants,
3.5 or 4.3\,d & ODE, CF & \nolinkurl{out/t55}:13--53, 111--153 \\
$k=10$ in place of $k=1$ & doubling delays shorter by under 1\,d; peaks 3\,d late in an
illness about 30\% higher & assumed dose, nominal constants & ODE &
\nolinkurl{out/t55}:13--20, 65--71, 95--104 \\
\end{longtable}
\endgroup
% src: as listed in the last column; REPORT.md:1167-1192 (R-T table; trough stricter in 6 of 8; k = 10 under a day, about 30%), :1194-1206 (R-U); out/t55_infusion_tolerance.txt:12-20,46,63-71,95-104,111-153,179; out/t30_monitoring_confounders.txt:13 (0.280 x 150 = 42), :23 (0.204 x 150 = 30.6); pi_note/pi_note.tex:395-396 (21 d capped at 4.3 d needs ~1.8%: 150.8 - 2850 exp(-21 ln2/4.3) = 54 uM/d)

% -----------------------------------------------------------------------------
\subsection{What the audit did not check}
\label{si:s8-limits}

All records are simulated from the models themselves, on constants that are partly
illustrative, so the analysis is a design study on a specified model and not an analysis of
a cohort; every absolute standard deviation depends on those constants. Verdicts within a few
per cent of the target of 0.20 depend on details. With $k_e$ and the load known exactly, 25
monthly draws gave 0.2044 under one pooled noise level and 0.1933 with a noise level for each
draw, and 0.198, 0.204 and 32.6 for draws spaced 29.2, 30.4 and 28 days apart, the last a
multiple of the dose interval. Where a structure published without a substrate load was given
one, the place it enters is our assumption, and three of the four choices tested leave
$\theta$ unidentifiable under the algebraic structure. The audit's literature figures (the
in vitro to in vivo mapping, the serial variability of propionylcarnitine, the breath-test
repeatability) come from other disorders, assays or healthy subjects and are used only to test
assumptions; none replaces a model value.
% src: REPORT.md:5 (design study on a specified, simulated model, not a cohort), :1369-1370 (0.2044 vs 0.1933; 0.198/0.204/32.6 at 29.2/30.4/28 d), :1417 (literature figures from other populations); paper/sections/A1_supplementary.tex:96-106 (load placement: three of four choices leave theta unidentifiable)
